\documentclass[12pt,a4paper]{article}
\usepackage{graphicx} % Required for inserting images
\usepackage{listings}     % для lstlisting
\usepackage{xcolor}       % для цветов в листингах
\usepackage[utf8]{inputenc}    % если pdflatex
\usepackage[T2A]{fontenc}      % если pdflatex
\usepackage[russian]{babel}
\usepackage{hyperref}% для русского текста
\usepackage{minted}
\usemintedstyle{friendly}
\title{Python}
\begin{document}

\begin{center}
\Large \textbf{Пособие для изучения Python} \\
\normalsize для подготовки к ОГЭ и ЕГЭ от \href{https://t.me/kvinnen}{Алины Лисянской}
\end{center}

\tableofcontents
% ==================== ПРЕДИСЛОВИЕ ====================
\section*{Предисловие}
\addcontentsline{toc}{section}{Предисловие}

Это пособие для повторение теории в Python. Оно поможет вам систематизировать знания, повторить синтаксис и освоить приёмы работы в Python, которые ускоряют написание кода на экзамене

% ==================== ВВЕДЕНИЕ ====================
\subsection*{Введение}
\addcontentsline{toc}{section}{Введение в Python}

Языки программирования бывают интерпретируемыми и компилируемыми. Компилируемые сначала переводят всю программу целиком в машинный код, а потом запускают. Интерпретируемые выполняют программу строчка за строчкой прямо во время работы. Python --- интерпретируемый. Отсюда вытекают его особенности: он медленнее многих других языков, но прост в изучении и имеет множество применений!

\section*{Переменные}
\addcontentsline{toc}{subsection}{Переменные}

Переменная -- это именованная ячейка памяти, в которой хранится какое-то значение. Создать переменную просто: пишем имя, знак равно и значение:

\begin{minted}{python}
x = 5
name = "Аня"
\end{minted}

Имя переменной может быть любым, но оно не может начинаться с цифры и не может содержать пробелы. Мы будем стараться называть переменные так, чтобы они отражали значение, которое хранят.
\medskip

\subsection*{Типы данных}
\addcontentsline{toc}{subsection}{Типы данных}

У каждого значения есть тип. Основные типы в Python:
\begin{itemize}
    \item \texttt{int} --- целое число;
    \item \texttt{float} --- дробное число;
    \item \texttt{str} --- строка, то есть текст;
    \item \texttt{bool} --- логическое значение (\texttt{True} или \texttt{False}).
\end{itemize}
\smallskip

Python сам определяет тип, когда мы создаём переменную:
\begin{minted}{python}
a = 10        # int
b = 3.14      # float
c = "Привет"  # str
\end{minted}
\smallskip

Иногда нам потребуется преобразовать один тип в другой. Для этого используют \texttt{int()}, \texttt{float()}, \texttt{str()}:

\begin{minted}{python}
int("5")   # строку "5" -> число 5
str(42)    # число 42 -> строку "42"
\end{minted}

\subsection*{Основные команды}
\addcontentsline{toc}{subsection}{Команды print и input}

\texttt{print()} выводит что-либо на экран. Можно передать текст в кавычках, число, посчитать арифметическое выражение, вывести значение переменной или сразу нескольких через запятую (в таком случае они выведутся через пробел):

\begin{minted}{python}
print("Привет!")           # Привет!
print(2 + 2)               # 4
print("Сумма:", 2 + 2)     # Сумма: 4
print(a, b, c)             # 10 3.14 Привет
\end{minted}

\textbf{input()} считывает то, что пользователь вводит с клавиатуры. Важно: \texttt{input()} всегда возвращает строку, даже если пользователь ввёл число. Поэтому если нужно, чтобы пользователь ввел не текст, а число, указываем нужный тип данных перед \texttt{input()}:

\begin{minted}{python}
name   = input()   # str
age    = int(input())  # int
weight = float(input())   # float
\end{minted}

\subsection*{Арифметические операции}
\addcontentsline{toc}{subsection}{Арифметические операции}

Познакомимся с основными арифметическими операциями. Все кроме деления здесь тривиальны:

\begin{minted}{python}
a = 10
b = 3

a + b   #сложение -> 13
a - b   #вычитание -> 7
a * b   #умножение -> 30
a ** b  #возведение в степень -> 1000
\end{minted}

С делением интереснее:
\begin{minted}{python}
a / b   #дробное деление -> 3.3333...
a // b  #целочисленное деление -> 3
a % b   #остаток от деления -> 1
\end{minted}

Обратите внимание: обычное деление \texttt{/} всегда возвращает \texttt{float}, даже если числа делятся нацело:

\begin{minted}{python}
10 / 2   #5.0 (а не 5)!
10 // 2  #5  (целочисленное деление вернёт int)
\end{minted}

Операция остатка от деления \texttt{\%} очень часто используется в задачах, например чтобы проверить чётность числа:

\begin{minted}{python}
if n % 2 == 0:
    print("чётное")
else:
    print("нечётное")
\end{minted}

Для удобства в Python есть \textbf{сокращённые операторы присваивания}:

\begin{minted}{python}
x = 10
x += 5   # то же самое, что x = x + 5  -> 15
x -= 3   # то же самое, что x = x - 3  -> 12
x *= 2   # то же самое, что x = x * 2  -> 24
x //= 4  # то же самое, что x = x // 4 -> 6
x %= 4   # то же самое, что x = x % 4  -> 2
\end{minted}

\textbf{Порядок выполнения операций} такой же, как в математике.
\bigskip

Также помним \textbf{про работу со строками}. Строки можно складывать.
\begin{minted}{python}
s1 = 'pine'
s2 = 'apple'
print(s1 + s2)  #'pineapple'
\end{minted}

А также умножать на число, по сути дублируя строку указанное количество раз:
\begin{minted}{python}
t1 = 'hi'
print(t1 * 3)  #'hihihi'
print(5 * t1) #'hihihihihi'
\end{minted}

\subsection*{Перевод чисел: bin(), oct(), hex()}
\addcontentsline{toc}{subsection}{Перевод чисел в другие СС}

Это встроенные функции для перевода чисел в другие системы счисления, которые нужны для некоторых заданий в ОГЭ и ЕГЭ.
\medskip

\begin{itemize}
    \item \textbf{bin(n)} переводит число в двоичную систему
    \item \textbf{oct(n)} -- в восьмеричную
    \item \textbf{hex(n)} -- в шестнадцатеричную
\end{itemize}

\begin{minted}{python}
bin(10)   # '0b1010' - префикс 0b: двоичная
oct(8)    # '0o10'   - префикс 0o: восьмеричная
hex(255)  # '0xff'   - префикс 0x: шестнадцатеричная
\end{minted}
\medskip

Если нужно получить число без префикса, используем срез строки (про них подробнее чуть позже):

\begin{minted}{python}
bin(10)[2:]   # '1010'
hex(255)[2:]  # 'ff'
\end{minted}
\medskip

Обратный перевод (из любой системы в десятичную) делает функция \textbf{int()} с двумя аргументами: первый -- строка с числом, второй -- основание системы счисления:

\begin{minted}{python}
int('1010', 2)   # 10  (из двоичной)
int('ff', 16)    # 255 (из шестнадцатеричной)
int('17', 8)     # 15  (из восьмеричной)
\end{minted}

\section*{Условный оператор}
\addcontentsline{toc}{section}{Условный оператор}

Условный оператор позволяет программе принимать решения: выполнять один блок кода или другой в зависимости от условия. Основная конструкция:

\begin{minted}{python}
if условие:
    #команда 1 выполнится, если условие истинно
elif другое_условие:
    #команда 2 выполнится, если первое условие ложно, а это истинно
else:
    #команда 3 выполнится, если все условия выше ложны
\end{minted}

Блоки кода внутри \texttt{if} обязательно выделяются отступом -- это называют табуляцией, равной 4 пробела. Python определяет принадлежность строки к блоку именно по этому отступу.
\medskip

Пример:

\begin{minted}{python}
age = int(input("Введи возраст: "))

if age < 0:
    print("Неверный возраст")
elif age < 18:
    print("Несовершеннолетний")
else:
    print("Совершеннолетний")
\end{minted}

Ветки \texttt{elif} и \texttt{else} необязательны. Имеется ввиду, что можно написать просто \texttt{if} и без них.

\smallskip
\noindent В условиях используются операторы сравнения:

\begin{minted}{python}
a == b   # a равно b
a != b   # a не равно b
a > b    # a больше b
a < b    # a меньше b
a >= b   # a больше или равно b
a <= b   # a меньше или равно b
\end{minted}

Несколько условий можно объединять логическими операторами \texttt{and} (и), \texttt{or} (или), \texttt{not} (не):

\begin{minted}{python}
x = 15
if x > 10 and x < 20:
    print("x находится между 10 и 20")
\end{minted}

\section*{Циклы}
\addcontentsline{toc}{section}{Циклы}
\subsection*{Цикл while}
\addcontentsline{toc}{subsection}{Цикл while}

Цикл \texttt{while} выполняет блок кода до тех пор, пока условие истинно:

\begin{minted}{python}
while условие:
    # тело цикла
\end{minted}

Пример -- считаем от 1 до 5:

\begin{minted}{python}
i = 1
while i <= 5:
    print(i)
    i += 1   # увеличиваем i на 1, иначе цикл бесконечный
\end{minted}

Важно следить за тем, чтобы условие когда-нибудь стало ложным -- иначе программа зависнет в бесконечном цикле.

Команда \texttt{break} прерывает цикл досрочно, а \texttt{continue} пропускает текущую итерацию и переходит к следующей:

\begin{minted}{python}
i = 0
while True:
    i += 1
    if i % 2 == 0:
        continue   # пропускаем чётные
    if i > 9:
        break      # останавливаемся после 9
    print(i)       # выведет 1, 3, 5, 7, 9
\end{minted}


\subsection*{Цикл for}
\addcontentsline{toc}{subsection}{Цикл for}

Цикл \texttt{for} перебирает элементы какой-либо последовательности -- строки, списка, диапазона чисел. На каждой итерации переменная цикла принимает очередное значение из последовательности:

\begin{minted}{python}
for переменная in последовательность:
    # тело цикла
\end{minted}

Разберём на примере. Напишем \texttt{for i in range(5)} --- это значит: «пусть переменная \texttt{i} по очереди принимает значения 0 до 5 не включительно, и для каждого из них выполняй тело цикла»:

\begin{minted}{python}
for i in range(5):
    print(i)    # выведет 0, затем 1, затем 2, затем 3, затем 4
\end{minted}

\medskip
\noindent Функция \texttt{range()} генерирует диапазон чисел и принимает от одного до трёх аргументов:

\begin{minted}{python}
range(stop)           # от 0 до stop - 1
range(start, stop)    # от start до stop - 1
range(start, stop, step)  # от start до stop - 1 с шагом step
\end{minted}

Важно запомнить: правая граница \textbf{не включается}. То есть \texttt{range(3, 6)} даст числа 3, 4, 5.
\bigskip

Пример с тремя аргументами:

\begin{minted}{python}
for i in range(0, 10, 2):
    print(i)    # 0, 2, 4, 6, 8

for i in range(10, 0, -1):
    print(i)    # 10, 9, 8, ..., 1 - отрицательный шаг, идём в обратном порядке
\end{minted}

Отрицательный шаг особенно полезен, когда нужно идти от большего к меньшему. При этом правило то же: правая граница не включается, поэтому \texttt{range(10, 0, -1)} даст числа от 10 до 1, но не 0.
\medskip

Цикл \texttt{for} тоже поддерживает \texttt{break} и \texttt{continue}.

\end{document}
